একটি \emph{প্রথম-ক্রমের ভাষা} \emph{ধ্রুবক}, \emph{অপেক্ষক} ও
\emph{বিধেয়} সংকেত নিয়ে গঠিত। অপেক্ষক ও বিধেয় সংকেত নির্দিষ্ট
সংখ্যক আর্গুমেন্ট নেয়; ধ্রুবক সংকেত কোনো আর্গুমেন্ট নেয় না।
যেমন, \emph{পাটিগণিতের ভাষায়} রয়েছে
একটি ধ্রুবক সংকেত~$\Obj 0$, একটি এক-স্থানীয় অপেক্ষক সংকেত $\prime$,
দুটি দ্বি-স্থানীয় অপেক্ষক সংকেত~$+$ ও $\times$, এবং একটি
দ্বি-স্থানীয় বিধেয় সংকেত~$<$। \emph{চলরাশি}, ধ্রুবক সংকেত ও
অপেক্ষক সংকেত থেকে আমরা ভাষার \emph{পদ} গঠন করি। ভাষার পদ,
বিধেয় সংকেত এবং \emph{অভিন্নতার সংকেত}~$\eq$ থেকে গঠন করি
\emph{পরমাণু সূত্র}। তারপর সেগুলি থেকে যৌক্তিক সংযোজক
$\lnot$, $\lor$, $\land$, $\lif$, $\liff$ এবং পরিমাণসূচক
$\lforall$ ও $\lexists$ ব্যবহার করে ভাষার সূত্র গঠন করি। পদ ও সূত্র
গঠনের সময় সব প্রয়োজনীয় বন্ধনী যথাস্থানে বসানোর ফলে কোনো সূত্র
পড়ার ঠিক একটিই উপায় থাকে। এর ফলে সূত্রের গঠনের উপর আরোহ করে
সংজ্ঞা দেওয়া সম্ভব হয়।

সূত্রে চলরাশির সংঘটন কখনো কখনো সংশ্লিষ্ট কোনো পরিমাণসূচকের অধীন
থাকে। কোনো চলরাশি সেই পরিমাণসূচকের \emph{পরিসরে} থাকলে তাকে
\emph{বদ্ধ} ধরা হয়, অন্যথায় \emph{মুক্ত}। এই ধারণাগুলির
সংজ্ঞা আরোহী। একটি সূত্রে কোনো চলরাশির জায়গায় একটি পদের
\emph{প্রতিস্থাপন} ক্রিয়াটিও আমরা আরোহীভাবে সংজ্ঞায়িত করি।
মুক্ত চলরাশির সংঘটন নেই এমন সূত্রকে \emph{বাক্য} বলে।
